\documentclass{article}
\usepackage{amsmath,hyperref}
\title{Fully observed OPEN does not certify an Open-RMF door will hold}
\author{}
\date{}
\begin{document}
\maketitle
\paragraph{Source boundary.}
The earlier \href{https://github.com/eternal-november/eos-research-programmes/pull/507}{door replay}
separates OPEN from a complete read of the expected leaves and tests
whether the fleet phase accepts an older OPEN with a new session.
Here every expected leaf is found and is sampled at its target in the
current update; the session matches. The pinned simulator's
\texttt{PreUpdate} calls \texttt{command\_door} before computing and
publishing its position-based mode. Thus the same complete, fresh
OPEN report can accompany a zero or closing write in the same update.
The source extraction executes the selected command, target-velocity,
mode and publication functions with inert entities. This does not run
\texttt{PreUpdate}, ROS transport, or Gazebo dynamics.

\paragraph{A currently OPEN target is not sufficient.}
After multiple closing writes, use an OPEN target in both histories,
holding both observed positions fixed at their exact open target
while the commands accumulate. The original limiter retains its last
command, not the measured leaf velocity. With timestep $1/8$ and
nominal acceleration $1/8$, let $a=1/64$ and let $k$ be the number of
prior closing ticks. At this read cut the negative closing magnitude is
$q=a\min(k,16)$, since the velocity cap is $1/4$ and the deceleration
condition for the selected closing distance does not bind. On each
OPEN update at the exact target, the limiter writes
\[
  u_{j}=-\max(q-ja,0).
\]
The initialized stable history and the residual-closing history have the
same position sample, OPEN target, report and fleet-phase decision.
The receiver restricted to these observations cannot justify a
persistent-open obligation in both. Merely checking that all leaves
were read and are currently open cannot repair it.

\paragraph{Separate modeled physical continuation.}
If each sampled position follows the issued velocity immediately
for one $1/8$-second step, the remaining closing displacement until
the first zero write is $m(m-1)/1024$, where $m=\min(k,16)$.
Before the zero write a displaced leaf remains below its open target,
and the nonpositive acceleration cap wins even outside the deadband,
so the formula does not assume the leaf stays inside it. In the
chosen tolerance $1/64$, four residual ticks in the initial limiter
state give a dip $12/1024$ and remain OPEN; five give $20/1024$
and do not. The branch is a conditional kinematic example, not a
measured Gazebo transition or a verified physical door-clearance
threshold. A continuing OPEN target also requires the authority to
exclude later close writes and a physical model capable of carrying
the held aperture forward through actual passage. A zero command is
not a measurement that the leaf has stopped.

\paragraph{Consequence for the shared receipt.}
The producer of the clearance report must account for remaining
effects as well as fresh observation and the current session.
A consumer with only the completed phase result does not see
whether the writer actually installed a safe continuation.
Giving a direct scheduler the same causal writer and plant evidence
would remove this distinction; the result is a source boundary,
not an advantage against an equally informed scheduler or an
independent mission test. Four pinned upstream files are specified
in the replay source note. The 420 frozen-position cases and 136
modeled moving update steps do not imply rates of occurrence.
\end{document}
